\section{Regulatory, Ethical, and Oversight Considerations}

The oversight architecture for this multicenter randomized controlled study layers
a conventional human-subject-protection structure with the additional governance
that an autonomy-bearing surgical system, deployed as a harmonized fleet across
eight sites, requires. A central Coordinating Center directs the conduct of the
trial across the network, and the Sponsor-Investigator, ChemicalQDevice, answers to
a single Institutional Review Board (sIRB) of record under 45 CFR 46.114
\cite{commonrule_sirb} and to an independent data and safety monitoring board
(DSMB) with a group-sequential interim plan. The structure convenes a Physical AI
Safety Review Committee on a $\leq 90$-day cadence, designates an Independent Safety
Monitor (ISM) for each procedure, and delegates day-to-day monitoring to a
robotics- and artificial-intelligence-competent clinical research organization
(CRO) over the site Principal Investigators and their teams. Figure 19 presents the
multicenter governance and safety-oversight structure.

\begin{mermaidfig}
\node[mmgoal](si)  at (4.2,0)     {Sponsor-Investigator\\ChemicalQDevice};
\node[mmstep](cc)  at (4.2,-2.0)  {Coordinating Center\\(8-site network)};
\node[mmstep](irb) at (0,-4)    {Single IRB (sIRB)\\45 CFR 46.114};
\node[mmstep](dsmb)at (4.2,-4.2)  {DSMB\\group-sequential\\interim};
\node[mmstep](pais)at (8.4,-4.2)  {Physical AI Safety\\Review Committee\\$\leq 90$-day cadence};
\node[mmin]  (cro) at (-2,-6.6)    {CRO / clinical\\monitor (robotics\\+ AI competent)};
\node[mmstep](ism) at (8.4,-6.6)  {Independent Safety\\Monitor (per\\procedure)};
\node[mmin]  (site)at (4.2,-8.4)  {Site PIs + teams:\\surgeon, operator,\\backup operator};
\draw[mmedge] (si) -- (cc);
\draw[->] (cc.south) -- (irb.north east);
\draw[mmedge] (cc) -- (dsmb);
\draw[->] (cc.south) -- (pais.north west);
\draw[mmedge] (cc.west) -| (cro.north);
\draw[mmedge] (cro) |- (site.west);
\draw[mmedged] (dsmb) -- node[right,font=\scriptsize]{stop/continue} (site);
\draw[mmedged] (pais.south) -- node[right,font=\scriptsize,pos=0.5]{USL changes} (ism.north);
\draw[mmedged] (ism) |- (site.east);
\end{mermaidfig}
\vspace{-0.65cm}
\figcaption{Figure 19. Multicenter governance and safety oversight. The
Sponsor-Investigator directs the trial through a Coordinating Center across the
eight-site network and answers to a single IRB and an independent DSMB with a
group-sequential interim plan, convenes a Physical AI SRC, designates an Independent Safety Monitor for each procedure,
and delegates monitoring to a robotics- and AI-competent CRO over site
Principal Investigators, teams.}

\subsection{Informed Consent Process}

Informed consent is obtained by a qualified member of the research team, who is not
the treating surgeon, before any study-specific procedure, in language the
participant can understand and with adequate time and opportunity to ask questions,
consistent with 21 CFR part 50 \cite{ecfr2026part50} and ICH E6(R3). A single
master consent form is used across the network under the single IRB. The consent
discussion covers the diagnosis, the alternatives, and the prognosis; the
randomized 1:1 allocation between Arm A (daraxonrasib at the recommended Phase 2
dose plus the LLM-directed robotic Whipple) and Arm B (modified FOLFIRINOX plus
standard high-volume pancreaticoduodenectomy) and the open-label nature of that
allocation; the risks of each regimen; and, in addition to the standard elements,
an explicit and distinct disclosure of the on-premises LLM-directed robotic system,
its role and continuous human oversight, and its system-specific risks, as required
for an autonomy-bearing investigational device. The consent form addresses the
therapeutic misconception directly and explains the staged-autonomy model under
which the system operates only at the supervised semiautonomous level.

A defining feature of this trial class is the documented Physical AI opt-out. Under
21 CFR $\S$312.60(f), as adapted for Physical AI investigations
\cite{kawchak2026adapt312}, a participant randomized to Arm A is offered, and may
exercise, a documented right to request administration without the Physical AI
advisory layer where clinically feasible, without penalty and without affecting
eligibility for the remainder of the study or for clinical care. The participant's
election or declination of the opt-out is recorded in the source document and the
CRF. A legally authorized representative process and the use of qualified
interpreters and translated materials for participants with limited English
proficiency are provided so that consent is meaningful for every eligible individual
across the network. The right to withdraw at any time without penalty is stated at
consent and reaffirmed at each contact. Figure 20 shows consent process and the
opt-out.

\begin{mermaidfig}
\node[mmin]  (intro)at (-0.2,0)      {Investigator presents\\study: disease,\\randomization, prognosis};
\node[mmstep](pai)  at (4.4,0)    {Explain Physical AI\\system: role, oversight,\\risks (\S312.60(f))};
\node[mmdec] (opt)  at (9.3,0)    {Participant\\choice};
\node[mmgoal](yes)  at (9.3,-2.6) {Consents to Physical AI\\signed ICF + assent\\if applicable};
\node[mmstep](no)   at (4.4,-2.6) {Requests non-Physical AI\\administration if\\clinically feasible};
\node[mmgoal](enr)  at (4.4,-4.9) {Randomize and document\\consent in source + CRF};
\node[mmin]  (wd)   at (-0.2,-4.9)   {Right to withdraw at\\any time without penalty};
\draw[mmedge] (intro) -- (pai);
\draw[mmedge] (pai) -- (opt);
\draw[mmedge] (opt) -- node[right,font=\scriptsize]{accept} (yes);
\draw[mmedge] (opt.south west) -- node[above,font=\scriptsize, xshift=-0.41cm, yshift=0.24cm]{decline AI} (no.east);
\draw[mmedge] (yes) |- (enr.east);
\draw[mmedge] (no) -- (enr);
\draw[mmedged] (enr) -- (wd);
\end{mermaidfig}
\vspace{-0.7cm}
\figcaption{Figure 20. Informed-consent process with the Physical AI opt-out.
Consent discloses the randomized\\allocation and the Physical AI system's role,
continuous human oversight, and system-specific risks,\\and offers a participant
randomized to Arm A a documented right to request non-Physical AI\\administration where clinically feasible (21 CFR $\S$312.60(f)), with the standard right to
withdraw.}

\subsection{Funding, Co-Investment Governance, and the Capital Firewall}

This Phase 2 trial is funded through a Patient-Aligned Co-Investment Facility
(PACIF), a mechanism by which wealthier individuals connected to cancer patients may
contribute ring-fenced philanthropic and private capital with the sole and explicit
purpose of raising the likelihood of trial success through operational capacity. The
facility purchases operational capacity only, and never scientific outcomes: capital
is directed exclusively to six operational levers, namely (1) the eight-site
academic HPB network, (2) the expanded Phase 0 simulation compute required to meet
the upgraded readiness gate, (3) the harmonized robot fleet across sites, (4) the
Patient Access and Equity Fund, (5) the central review and biomarker core, and (6)
the independent oversight bodies. None of these levers touches the science of the
comparison.

The integrity of the comparison is protected by a capital firewall, a structural
separation that prevents any funder or contributor from influencing the
randomization, the endpoint definitions, the outcome adjudication, the statistical
analysis, the data access, or the publication of the trial. Contributions enter
through an independent fund administrator and are disbursed only against the
permitted operational levers; no funder holds any role, vote, or veto over the
scientific conduct, and no contribution is conditioned on any result. All financial
interests are disclosed and managed under 21 CFR part 54 \cite{ecfr2026part54}, and
the facility is governed by the verification-and-validation-and-uncertainty-quantification
(VVUQ) financial-data standard of the H.R. 9510 framework \cite{kawchak2026hr9510},
which ties the funding and oversight of Physical AI oncology investigations to a
verifiable evidentiary standard. The Patient Access and Equity Fund is administered
for equitable disbursement, supporting participation costs and access across the
network, and is itself walled off from eligibility determinations and outcome
ascertainment, so that no participant's access to the trial or to its assessments
can be influenced by the source of funds. All results, regardless of outcome, are
openly deposited under the Creative Commons Attribution 4.0 International (CC BY 4.0)
license. Figure 21 shows the capital firewall governance.

\begin{mermaidfig}
\node[mmin]  (fund) at (0,0)        {Funders / PACIF\\(patient-aligned\\private + philanthropic)};
\node[mmdec] (fw)   at (0,-3.0)     {Capital\\firewall};
\node[mmstep](admin)at (0,-5.6)     {Independent fund\\administrator\\(disburse-only)};
\node[mmstep](lev)  at (-4.6,-3.0)  {Permitted operational\\levers: 8 sites,\\Phase 0 compute, fleet,\\access fund, review core};
\node[mmgoal](ovr)  at (0,-7.82)     {Independent oversight\\sIRB, DSMB, Physical\\AI Safety Committee};
\node[mmdark](sci)  at (5.3,-3.0)   {Randomization,\\endpoints, adjudication,\\analysis, data,\\publication};
\draw[mmedge]  (fund) -- (fw);
\draw[mmedge]  (fw) -- (admin);
\draw[mmedge]  (admin.north west) -- (lev.south east);
\draw[mmedge]  (admin) -- (ovr);
\draw[mmedged] (fw.east) -- node[above,font=\scriptsize,pos=0.5, yshift=-0.45cm, xshift=0.0cm]{blocked} (sci.west);
\node[font=\scriptsize\bfseries,text=protomaroon] at (3.65,-1.5) {no influence};
\end{mermaidfig}
\vspace{-0.7cm}
\figcaption{Figure 21. Capital firewall governance. Patient-aligned funders
contribute through the PACIF into an independent fund administrator behind a capital
firewall that disburses only to the permitted operational levers (the eight-site
network, Phase 0 compute, the robot fleet, the Patient Access and Equity Fund, and
the central review and biomarker core) and to independent oversight; the dashed
blocked path shows that funders cannot reach randomization, endpoint definition,
outcome adjudication, statistical analysis, data access, or publication.}

\subsection{Study Discontinuation and Closure}

The study may be discontinued or closed by the Sponsor-Investigator, by the single
IRB, or by the FDA. Grounds for closure include unacceptable safety findings
(including the prespecified DSMB stopping rules of $\S$9), a clinical hold, a
determination that the study cannot be conducted in compliance, or completion of the
planned accrual and analysis. The FDA may impose a clinical hold on the drug or
device component under the provisions for terminating or suspending an
investigation; in addition, the Physical AI overlay establishes specific
clinical-hold grounds for the autonomy-bearing system, including a breach of a hard
cyber-physical safety limit, a failure of the verification-before-generation gate
surface, or a loss of audit integrity, and these grounds may halt activity at a
single site or across the fleet. On any hold or closure, enrollment and
patient-facing robotic activity stop immediately at the affected scope; affected
participants are transitioned to appropriate clinical care; and a controlled
decommissioning and data-preservation procedure is executed, under which the robotic
system is locked out of patient contact, the deterministic configuration and
repository commit are recorded, and the complete federated hash-chained command and
telemetry record is sealed and retained for the regulatory retention period
($\S$10.10). The single IRB, the FDA, and the DSMB are notified, and a final report
is prepared.

\subsection{Confidentiality and Privacy}

Participant confidentiality is protected throughout. Data released for analysis,
deposition, or publication are de-identified by the HIPAA Safe Harbor method, with
removal of all eighteen categories of protected health information identifiers
(names; geographic subdivisions smaller than a state; all date elements other than
year that are directly related to an individual, and ages over 89; telephone and
fax numbers; email addresses; Social Security numbers; medical record numbers;
health plan beneficiary numbers; account numbers; certificate or license numbers;
vehicle identifiers; device identifiers and serial numbers; web URLs; IP addresses;
biometric identifiers; full-face photographs and comparable images; and any other
unique identifying number, characteristic, or code). Access to identifiable records
is governed by a deny-by-default authorization model, in which no role can read,
write, or export a record except where an access right has been explicitly granted,
and every access, change, and export is written to a federated hash-chained audit
trail in which each entry cryptographically incorporates the hash of the prior entry
so that the record is tamper-evident and any alteration is detectable across the
eight sites. All electronic records, signatures, and the cyber-physical telemetry
stream are maintained under 21 CFR part 11 \cite{ecfr2026part11}, with validated
systems, secure time-stamped audit trails that do not obscure prior entries,
role-based access controls, and electronic signatures bound to the records they
authenticate.

\subsection{Future Use of Stored Specimens and Data}

The cyber-physical record is retained as a tiered telemetry pyramid that balances
long-term scientific reuse against storage cost. At the apex are the compact,
human-readable summaries and the safety-event and audit logs retained in full and
indefinitely; the middle tiers hold compressed feature and event streams; and the
base, level zero (L0), holds the raw, losslessly reconstructable command and sensor
stream. The L0 raw tier is deposited on Zenodo under controlled terms and is
decompressible and reconstructable on FDA request as required for inspection and
review under 21 CFR $\S$312.57, so that any reported behavior can be replayed
deterministically from the seed and the repository commit. Circulating tumor DNA
(ctDNA) and other biomarker specimens collected for the central biomarker core are
governed by the participant's consent and applicable regulation, are processed for
the protocol-specified assessments, and are not retained for an unspecified future
purpose without appropriate consent. Any future secondary use of stored
de-identified data is limited to research consistent with the consent, is conducted
on de-identified data only, and preserves the audit chain.

\subsection{Key Roles and Study Governance}

The Sponsor-Investigator (ChemicalQDevice, Kevin Kawchak) holds the combined
IND/IDE and bears the sponsor and investigator responsibilities under 21 CFR parts
312 and 812, including protocol custody, safety reporting, oversight of the
investigation, and assurance of compliance across the network. The Coordinating
Center provides central operational direction, the single IRB of record, the
central randomization and review services, and harmonization across sites. Each site
Principal Investigator directs the conduct of the study at that site and supervises
the research team. At each procedure in Arm A, the operating surgeon performs the
pancreaticoduodenectomy under the supervised semiautonomy model; a primary operator
drives the system, and a qualified backup operator is present and ready to assume
control or to invoke the manual override at any time. The CRO provides robotics- and
AI-competent clinical monitoring across sites. The DSMB, the Independent Safety
Monitor, and the Physical AI Safety Review Committee provide the independent
oversight described in $\S$10.9. Roles, delegated duties, and the corresponding
qualifications are documented on a delegation-of-authority log maintained at each
site, and changes are recorded and reported as required.

\subsection{Safety Oversight}

Independent safety oversight operates on three tiers. The DSMB, composed of members
independent of the conduct of the study with surgical, oncologic, biostatistical,
and device-safety expertise, reviews accumulating safety and efficacy data,
executes the prespecified group-sequential interim analysis with O'Brien-Fleming
spending \cite{DeMets1994lan}, applies the prespecified stopping rules of $\S$9, and
recommends continuation, modification, pause, or halt for the trial or for an
individual site. The Independent Safety Monitor is designated for each procedure and
provides real-time, case-level safety oversight independent of the operating team,
with the authority to call a stop. The Physical AI Safety Review Committee meets on a
$\leq 90$-day cadence (and ad hoc on any triggering event) to review the
autonomy-bearing behavior of the fleet: the telemetry and audit record, the
verification-before-generation gate-surface decisions, the Unified Safety Level
(USL) rating, and any software change, and to recommend USL reassessment or
oversight changes to the Sponsor-Investigator. The three tiers are coordinated so
that no autonomy-bearing change and no escalation decision proceeds without
appropriate independent review.

\subsection{Clinical Monitoring}

The study is monitored by a CRO whose monitors are competent in both clinical-trial
monitoring and the robotics and artificial-intelligence elements of this trial
class. Monitoring follows a risk-based plan calibrated to site volume and risk, and
includes source-data verification not only of conventional clinical source
documents and the consent records but also of the cyber-physical evidence at each
site: the device telemetry, the verification-before-generation gate-surface logs,
and the federated hash-chained audit trail, which are verified against the recorded
outcomes and the deterministic seed and repository commit. Monitoring visits verify
eligibility, randomization integrity, consent (including the documented Physical AI
opt-out election in Arm A), safety reporting, protocol adherence, and data
integrity across the network, and findings are documented and resolved through to
closure.

\subsection{Quality Assurance and Quality Control}

Quality assurance and quality control span the clinical and the cyber-physical
domains and are harmonized across the eight sites. Before every procedure in Arm A,
the pre-procedure safety matrix is verified, the USL readiness rating is confirmed at
or above the deployment threshold, and the emergency-stop, positional, and force
envelopes are checked. After any software update, model change, or configuration
change to the LLM-directed system, a Phase 0 re-validation is performed against the
independent simulation frameworks, the USL is reassessed against the upgraded
threshold of at least 8.0, and fleet harmonization is reconfirmed so that every
device instance behaves identically, before any further patient-facing use at any
site; no update reaches a participant without re-validation. Audits may be conducted
by the Sponsor-Investigator or a designee independent of the conduct of the study,
and corrective and preventive actions are documented. The deterministic,
version-controlled build ensures that the validated configuration in use is exactly
the configuration that was tested and recorded.

\subsection{\hspace{0.03cm} Data Handling and Record Keeping}

Seven Physical AI record types are captured and retained for this trial class, in
addition to the conventional clinical trial records: (1) the LLM advisory record
(each prompt, context, and advisory output); (2) the robot command record (each
command issued to the arms); (3) the sensor and telemetry record (the
cyber-physical stream, tiered as in $\S$10.5); (4) the safety-limit-event record
(emergency-stop activations and latencies, positional and force-envelope
excursions, and vascular-exclusion gating events); (5) the
verification-before-generation gate-surface record (each accept, block, or escalate
decision); (6) the human-override and intervention record; and (7) the federated
hash-chained audit trail that binds the preceding six to a deterministic seed and a
repository commit across the network. All records are retained for the period
required by 21 CFR $\S$312.57 (at least two years after a marketing application is
approved for the indication, or, if no application is filed or approved, two years
after the investigation is discontinued and the FDA is notified), and the L0 raw
tier is preserved as described in $\S$10.5 so that a requested record can be
reconstructed.

\subsection{\hspace{0.03cm} Protocol Deviations}

A protocol deviation is any departure from the IRB-approved protocol. Deviations are
identified, documented, assessed for their effect on participant safety and data
integrity, and reported to the single IRB and the Sponsor-Investigator in accordance
with the IRB's policy and applicable regulation; deviations that affect safety or
the rights and welfare of participants, and any that recur, are escalated, and the
Coordinating Center tracks deviation patterns across sites. Deviations are
distinguished from planned protocol amendments, which follow the amendment process.
Where a deviation is necessary to eliminate an apparent immediate hazard to a
participant, the change may be implemented without prior IRB or FDA approval under 21
CFR $\S$312.30(b) (and the corresponding device provision), with prompt notification
and reporting afterward; for the autonomy-bearing system, an immediate-hazard
modification also triggers the decommissioning and re-validation controls of
$\S$10.4 and $\S$10.11 before any further patient-facing use.

\subsection{\hspace{0.03cm} Publication and Data Sharing Policy}

Results are disseminated regardless of outcome through dual deposition: the analysis
code, derived data sets, and protocol are published in a public GitHub repository
and archived on Zenodo with a citable DOI, under the CC BY 4.0 license. Reporting
follows the TRIPOD+AI statement for the model-assisted advisory component
\cite{Collins2024tripodai} and the CONSORT statement for the parallel-group
randomized design \cite{Schulz2010consort}, in addition to the applicable
clinical-trial reporting standards. The deterministic seed and the repository commit
are published so that every reported result is reproducible, and the tiered
telemetry record is deposited as described in $\S$10.5. The results of the
co-investment-funded trial are openly deposited so that the funding model is
transparent and the science is independently verifiable. No identifiable participant
information is published; all shared data are de-identified by the HIPAA Safe Harbor
method ($\S$10.5).

\subsection{\hspace{0.03cm} Conflict of Interest Policy}

Financial interests are disclosed and managed in accordance with 21 CFR part 54
(Financial Disclosure by Clinical Investigators) \cite{ecfr2026part54}. The
Sponsor-Investigator is the chief executive of ChemicalQDevice, which develops the
Physical AI system under study; this relationship is disclosed in full, and the
independent oversight structure of $\S$10.9 (the DSMB, the Independent Safety
Monitor, and the Physical AI Safety Review Committee) is established in part to
manage it, so that safety and escalation decisions are not made by a party with a
financial interest in the outcome alone. The co-investment capital is managed behind
the capital firewall described in $\S$10.3, so that no funder, including any with a
financial interest, can influence the randomization, the endpoints, the
adjudication, the analysis, the data access, or the publication. Investigator
financial-disclosure certifications are collected before participation and updated
through one year after the study, and any identified conflict is disclosed and
managed before it can affect the conduct, analysis, or reporting of the study.
